\chapter{Honesty, Axioms, and the Art of Trusting Proofs}
\label{ch:honesty}

\section{``Formally verified'' is a claim, not a spell}

The phrase \emph{formally verified} has marketing gravity, and marketing
gravity attracts abuse. This chapter arms you with the auditor's toolkit:
what a Lean certificate actually rests on, how to interrogate it in one
command, and the specific ways a verification claim can be hollow while
every file compiles. Nothing here is hypothetical --- each failure mode
appears in the wild, and the discipline below is the one the companion
projects hold themselves to.

\section{The trust ledger}

When the kernel accepts \lean{fieldImplementation}, what exactly are you
being asked to believe? Lean can tell you --- precisely:

\begin{lstlisting}[language=Lean]
#print axioms fieldImplementation
-- 'fieldImplementation' depends on axioms:
-- [propext, Classical.choice, Quot.sound]
\end{lstlisting}

\lean{\#print axioms} walks the \emph{entire} dependency tree of a theorem
--- every lemma, every lemma's lemmas, down to bedrock --- and reports every
assumption found there. The three names above are Lean's standard trio
(propositional extensionality, classical choice, quotient soundness):
ordinary classical mathematics, accepted by working mathematicians,
scrutinized by logicians for a century. A certificate reporting exactly
these three is called \textbf{axiom-clean}. Anything \emph{else} in that
list is a custom assumption someone added --- and the whole audit consists
of reading that list and asking whether you believe each entry.

\begin{bigidea}
A machine-checked theorem is a receipt with a complete list of its own
assumptions --- something prose mathematics has never had. But the receipt
only protects people who read it. \textbf{The one-command audit:} run
\lean{\#print axioms} on the headline theorem; anything beyond
\lean{[propext, Classical.choice, Quot.sound]} is where the bodies are
buried. Make this reflex, and no verification claim can hide from you.
\end{bigidea}

\section{A field guide to hollow certificates}

Compiling proofs can still be worthless. The four classic ways, in
increasing order of subtlety:

\textbf{1. The \lean{sorry}.} Lean's placeholder accepts any goal with a
warning. Fine for work in progress; fatal in a certificate, because
downstream theorems inherit the hole silently. Detection: the compiler
warns, and \lean{\#print axioms} shows \lean{sorryAx}. Trivial to catch ---
if you look.

\textbf{2. The smuggled axiom.} Stuck on a lemma? \lean{axiom} makes it
true. Sometimes legitimate (see the trusted-base section below); rotten when
undisclosed --- a proof ``of'' the group law that axiomatizes the hard
half of the group law is theater. Detection: \lean{\#print axioms}, always.

\textbf{3. The trivial specification.} The most instructive one. Consider:

\begin{lstlisting}[language=Lean]
theorem mul_correct : ∀ a b, ∃ c, mul a b = c := by
  intro a b; exact ⟨_, rfl⟩   -- checks! and says NOTHING
\end{lstlisting}

Kernel-approved, axiom-clean, and utterly empty: ``mul returns whatever it
returns.'' No tool catches this, because nothing is wrong \emph{formally}
--- the defect is that the statement doesn't say what the reader assumes it
says. The only detector is a human reading the \emph{statement} (never mind
the proof) and asking: \emph{if the code were wrong, would this theorem
fail?} For the trivial spec, the answer is no --- a buggy \code{mul}
satisfies it identically.

\textbf{4. The wrong model.} The proof is real, the spec is strong --- but
the thing verified isn't the thing that ships: a hand-transcription
(Chapter~\ref{ch:rust}), a stale extraction, a simplified semantics.
Detection: trace the chain from artifact to source --- extraction scripts,
pinned tool versions, regeneration instructions. If the chain can't be
replayed, the claim is about an orphan.

\begin{pitfall}
Rank these by danger and notice the inversion: the crude failures
(\lean{sorry}, smuggled axioms) are machine-detectable in seconds, while
the subtle ones (trivial specs, wrong models) defeat every automated check
and yield only to a thoughtful reader. Verification does not eliminate the
need for human judgment; it \emph{concentrates} all of it into two small,
well-lit places --- the statement and the model. That concentration is the
gift; squandering it by not reading the statement is the sin.
\end{pitfall}

\section{Honest boundaries: the trusted base}

Real projects meet real limits: SHA-512's compression function, SIMD
backends no translator models, foreign function calls. The honest move is
not to pretend, but to \emph{declare}: state the unverified piece as an
explicit assumption, document it in a ledger (the companion repos call
theirs \code{TRUSTED-BASE.md}), and let \lean{\#print axioms} carry the
disclosure to every downstream theorem automatically.

\begin{lstlisting}[language=Lean]
-- Declared, documented, and visible in every audit forever:
axiom sha512_spec : ∀ msg, Sha512.hash msg = SHA512_ideal msg
\end{lstlisting}

A signature-layer certificate honestly reads: \emph{EdDSA verification is
correct, GIVEN the hash behaves ideally and GIVEN the documented backend
assumptions} --- with both \emph{given}s machine-visible. Compare the two
postures: ``everything verified!'' (and hope nobody checks) versus ``these
two assumptions, this ledger, audit me'' --- the second is both humbler and
\emph{stronger}, because its claim survives the audit.

The companion projects add one more layer of candor worth copying: a
\emph{failure map}. Their control repository documents the dead ends ---
tactic patterns that exhaust memory, extraction scopes that drag in the
world, proof styles that don't scale --- each with the tell that identifies
it early. Knowledge of where the cliffs are is part of the method, and
pretending the cliffs don't exist is how the next person walks off one.

\begin{aha}
Notice the running theme: at every level, the methodology converts
\emph{invisible} trust into \emph{visible} trust. Extraction made the
code-to-model step visible; two-clause specs made operating envelopes
visible; \lean{\#print axioms} makes logical debts visible; the trusted-base
ledger makes engineering limits visible. Formal verification's deepest
product is not certainty --- it is \textbf{legibility of exactly what
remains uncertain}.
\end{aha}

\begin{tryit}
Audit the real thing. In \code{dalek-ed25519-verified}, run
\lean{\#print axioms} on \code{fieldImplementation} and
\code{edwardsImplementation} --- confirm the clean trio. Then read
\code{TRUSTED-BASE.md} and match each entry to where it would surface in an
audit. Finally, write a deliberately trivial spec for \code{add}, prove it
in one line, and observe that every automated check passes. Keep that file
open for one full minute. That minute is the chapter.
\end{tryit}

\section*{Exercises}

\exercise{For each hollow-certificate species, name its detector: (a)
\lean{sorry}; (b) smuggled axiom; (c) trivial spec; (d) wrong model. Which
two can a CI pipeline catch mechanically, and what CI check would you write
for each?}

\exercise{Strengthen this spec until a buggy implementation would fail it:
\lean{theorem sub_ok : ∀ a b, ∃ c, sub a b = .ok c}. (List what's missing:
bounds hypotheses? bounds propagation? the value equation? Compare with
Chapter~\ref{ch:denotation}'s two-clause shape.)}

\exercise{A vendor's whitepaper says: ``Our signature library is formally
verified in Lean.'' Draft the five questions you would send them, in
priority order, and the answer you would require for each before relying on
the claim. (You now know all five.)}

\exercise{(Discussion) The trusted-base \lean{axiom} for SHA-512 and the
smuggled \lean{axiom} for a hard lemma are the \emph{same language feature}.
Articulate the difference in one sentence --- it is not technical.}

\begin{checkpoint}
You should now be able to: run and interpret the one-command audit; name
the standard three axioms and greet anything else with suspicion; explain
why trivial specs and wrong models defeat automation and what defeats
\emph{them}; and argue --- with conviction --- why a declared trusted base
is stronger, not weaker, than a claim of totality.
\end{checkpoint}
